\section{Реализация алгоритмов}

\subsection{Быстрая сортировка на C++}

Реализация быстрой сортировки выполнена на языке C++ с использованием
стандартной библиотеки \texttt{<vector>} и \texttt{<chrono>} для замера
времени выполнения. Полный исходный код приведён в приложении~\ref{app:quicksort}.

Ключевой функцией является \mintinline{cpp}{partition}, которая разделяет
подмассив \mintinline{cpp}{arr[low..high]} относительно опорного элемента
\mintinline{cpp}{arr[high]} и возвращает его итоговый индекс.
Ниже приведён фрагмент основной рекурсивной функции:

\begin{minted}{cpp}
void quickSort(std::vector<int>& arr, int low, int high) {
    if (low < high) {
        int pi = partition(arr, low, high);
        quickSort(arr, low, pi - 1);
        quickSort(arr, pi + 1, high);
    }
}
\end{minted}

Функция \mintinline{cpp}{quickSort} принимает ссылку на вектор \mintinline{cpp}{arr},
а также границы текущего подмассива"--- \mintinline{cpp}{low} и \mintinline{cpp}{high}.
Рекурсия завершается, когда \mintinline{cpp}{low >= high}, то есть подмассив
содержит не более одного элемента~\cite{cppreference}.

\subsection{Сортировка слиянием на Python}

Реализация сортировки слиянием выполнена на языке Python~3. Полный исходный
код приведён в приложении~\ref{app:mergesort}.

Алгоритм состоит из двух функций. Функция \mintinline{python}{merge_sort}
рекурсивно делит массив пополам, а функция \mintinline{python}{merge}
выполняет слияние двух отсортированных подмассивов. Ниже приведён
фрагмент функции слияния:

\begin{minted}{python}
def merge(arr: list, left: int, mid: int, right: int) -> None:
    left_part  = arr[left:mid + 1]
    right_part = arr[mid + 1:right + 1]
    i = j = 0
    k = left
    while i < len(left_part) and j < len(right_part):
        if left_part[i] <= right_part[j]:
            arr[k] = left_part[i]; i += 1
        else:
            arr[k] = right_part[j]; j += 1
        k += 1
\end{minted}

Переменные \mintinline{python}{i}, \mintinline{python}{j} и
\mintinline{python}{k} являются индексами левого подмассива,
правого подмассива и результирующего массива соответственно.
Стабильность алгоритма обеспечивается строгим неравенством
\mintinline{python}{left_part[i] <= right_part[j]}~\cite{pythondocs}.
